\section{Hardness of Matroid Optimization with a Linear Constraint}
\label{sec:tight_absolute}
In this section we use \Cref{thm:emb} %and \Cref{cor:deter} 
to prove \Cref{thm:main}. We apply the following reductions from EMB to MOL problems. Recall that $\cQ = \cP \setminus \{(\min,\IS,\leq)\}$ is the set of parameters for non-trivial MOL problems. Given a $P$-MOL problem for some $P \in \cQ$, and an EMB instance $I$, the reduction returns an instance $R_P(I)$ of the $P$-MOL problem. Note that the reduction is purely mathematical, and 
% in the sense that it 
does not specify the encoding of the instance. This will be useful for obtaining our hardness results in non-oracle computational models (see \Cref{sec:SAT}). 

Given $(\opt,\cF,\triangleleft), (\opt',\cF',\triangleleft') \in \cQ$, we use the notation $(\opt \textnormal{ is } \opt')$,  $(\cF \textnormal{ is } \cF')$, $(\triangleleft \textnormal{ is } \triangleleft')$ to denote the boolean expressions of equality between parameters of a MOL problem. For example, $(\opt \textnormal{ is } \opt')$ is {\em true} if and only if either $\opt,\opt'$ are both $\max$, or $\opt,\opt'$ are both $\min$. 

%$(\opt \textnormal{ is } \max)$ and $(\opt' \textnormal{ is } \max)$ are true, or $(\opt \textnormal{ is } \min)$ and $(\opt' \textnormal{ is } \min)$ are true. 


\begin{definition}
	\label{def:reduction}
	Given an \textnormal{EMB} instance $I =(E,\cI,c,T)$ and  $P \in \cQ$ where $P = (\opt,\cF,\triangleleft)$, define the
	 {\em reduced $P$-MOL instance} of $I$, denoted by  $R_P(I) = (E,\cI,v_I,w_{I,P},L_{I,P})$, as follows.
	\begin{enumerate}
		\item Define the auxiliary variable 
		\begin{align*}
			%	n & : \text{The variable representing a natural number.} \\
			d(P) = & \begin{cases}
				0 & ~~\textnormal{if  } \bigg((\opt \textnormal{ is } \max)  \textnormal{ and }   (\triangleleft \textnormal{ is}   \leq) \bigg) \textnormal{ or } 
				\bigg((\opt \textnormal{ is } \min)  \textnormal{ and }   (\triangleleft \textnormal{ is}   \geq) \bigg)\\
				1 & ~~\textnormal{otherwise.}
			\end{cases}
		\end{align*} For example, if $P = (\max,\IS,\leq)$ then $d(P) = 0$, and if  $P' = (\max,\IS,\geq)$ then $d(P') = 1$. 
		\item Let $H_I = 2 \cdot \max \{1, c(E)\}$.
		\item For all $e \in E$ let $v_I(e) =  H_I +c(e)$. 
		\item For all $e \in E$ let $w_{I,P}(e) =  H_I +c(e) \cdot (-1)^{d(P)}$.
		\item Let $k_I = \max_{S \in \cI} |S|$ be the rank of $(E,\cI)$. 
		\item Define $L_{I,P} = k_I \cdot H_I +T \cdot  (-1)^{d(P)}$. 
	\end{enumerate}
\end{definition}


%learly, given an EMB instance $I$ and $P \in \cQ$ the $P$-MOL instance $R_P(I)$ can be computed in polynomial time. 
%The main properties of the reduction are summarized in the following theorem. 
Now, for every EMB instance $I =(E,\cI,c,T)$, define the {\em error parameter} of $I$ as \begin{equation}
\label{eq:epsI}
\eps_I =  \frac{1}{8 \cdot (|E|+1) \cdot (T+1) \cdot \left(c(E)+1\right)}.
\end{equation}
Indeed, since the value selected for of the error parameter is sufficiently small, we can use a $(1+\eps_I)$-approximation for $R_P(I)$ to decide an EMB instance $I$.  
\begin{restatable}{theorem}{thmH}
	\label{thm:H}
	Given an instance $I =(E,\cI,c,T)$ of \textnormal{EMB}, and $P \in \cQ$ with $P = (\opt,\cF,\triangleleft)$, the following holds. 
	\begin{enumerate}
		\item If there is a solution $S$ for $R_P(I)$ such that $v_I(S) = k_I \cdot H_I +T$, then $I$ is a ``yes"-instance.  \label{h:1}
		\item  If $I$ is a ``yes"-instance then: %$I$ has a solution then: 
		$(i)$ $R_P(I)$ has a solution, and $(ii)$ every $(1+\eps_I)$-approximate solution $S$ for $R_P(I)$ satisfies $v_I(S) = k_I \cdot H_I +T$. \label{h:2}
	\end{enumerate}
\end{restatable}
Using \Cref{thm:H} and an assumed randomized Fully PTAS for the $P$-MOL problem, we can decide EMB in time which contradicts \Cref{thm:emb}. This gives the proof of \Cref{thm:main}. We first prove \Cref{thm:main}, and later give the proof of \Cref{thm:H}.   

 \thmMain*

\begin{proof}
	Assume towards a contradiction that there is a randomized Fully PTAS $\cA$ for oracle $P$-MOL. We use $\cA$ to decide oracle-EMB. Let $I = (E,\cI,c,T)$  be an oracle-EMB instance, and %let $n = |E|+1$ and let $m = c(E)+1$; observe that $n,m \neq 0$. C
	consider the following randomized algorithm $\cB$ that decides $I$. 
	\begin{enumerate}
		\item Construct the oracle $P$-MOL instance $R_{P}(I)$
		with the membership oracle of $(E,\cI)$.\label{step:1} 
		\item 
		Execute $\cA$ with the input $R_{P}(I)$ and $\eps_I$. \label{step:2} 
		\item If $\cA$ returns that $R_{P}(I)$ does not have a solution $-$ Return ''no" on $I$. \label{step:3}  %does not have a solution. \label{step:3} 
		\item Otherwise, let $S \leftarrow \cA(R_{P}(I),\eps_I)$ be the solution returned  by $\cA$.\label{step:4} 
		\item Return ''yes" on $I$ if and only if $v_I(S) = H_I \cdot k_I+T$. \label{step:5} 
	\end{enumerate}
	
	Let $n = |E|+1$ and $m = c(E)+1$. Note that $R_{P}(I)$ can be naively constructed from $I$ in time $\left(n \cdot (T+2) \cdot m\right)^{O(1)}$ using \Cref{def:reduction}. 
	As $\cA$ is a randomized Fully PTAS for oracle $P$-MOL, and by the selection of the error parameter \eqref{eq:epsI}, the running time of $\cB$ on $I$ is $\left(n \cdot (T+2) \cdot m\right)^{O(1)}$. %For the correctness,
	We now show correctness. 
	\begin{itemize}
		\item If $\cB$ returns ``yes" on $I$, then Step~\ref{step:4} of the algorithm computes a solution $S$ for $R_p(I)$ satisfying $v_I(S) = H_I \cdot k_I+T$. Thus, by \Cref{thm:H}, $I$ is a ``yes" instance. %has a solution.   
		
		
		\item If $I$ is a ``yes" instance then $R_P(I)$ has a solution by \Cref{thm:H}.
		As $\cA$ is a randomized Fully PTAS, with probability at least $\frac{1}{2}$ $\cA$ returns a $(1+\eps_I)$-approximate solution $S$ for $R_P(I)$ in Step~\ref{step:4}. By \Cref{thm:H}, $v_I(S) = H_I \cdot k_I+T$ (with probability at least $\frac{1}{2}$). Thus, $\cB$ returns ``yes" on $I$ with probability at least $\frac{1}{2}$.
		\end{itemize}  Hence, $\cB$ is a randomized algorithm which decides the oracle-EMB instance $I$ in time \newline $\left(n \cdot (T+2) \cdot m\right)^{O(1)}$. This is a contradiction to \Cref{thm:emb}.  
\end{proof}













In the remainder of this section we prove \Cref{thm:H}. We start with some basic properties of the reduction outlined in \Cref{def:reduction}.   

\begin{lemma}
\label{lem:FD2}
Given an instance $I =(E,\cI,c,T)$ of \textnormal{EMB},
let $P \in \cQ$ and consider a solution $S$ for $R_P(I)$ satisfying $v_I(S) = k_I \cdot H_I+T$. Then, $S$ is a solution for $I$.  
\end{lemma}

\begin{proof}
	Let $\cm = (E,\cI)$. As $S$ is a solution for $R_P(I)$, it holds that $S \in \IS(\cm)$; thus, $|S| \leq k_I$. Assume towards contradiction that $|S|<k_I$. Then, 
	$$v_I(S) = |S| \cdot H_I+c(S) \leq (k_I-1) \cdot H_I+c(S) \leq (k_I-1) \cdot H_I+c(E)<k_I \cdot H_I \leq k_I \cdot H_I+T.$$ 
	We reach a contradiction since $v_I(S) = k_I \cdot H_I+T$; thus, $|S| = k_I$, and  
	\begin{equation}
		\label{eq:S=k}
		k_I \cdot H_I+T = v_I(S) =   |S| \cdot H_I+c(S) =  k_I \cdot H_I+c(S). 
	\end{equation}  As $|S| = k_I$, we have that $S$ is a basis of $\cm$, and by \eqref{eq:S=k}, $c(S) = T$. Hence, $S$ is a solution for $I$.   
\end{proof}

The next result is the converse of the statement in \Cref{lem:FD2}.   

	\begin{lemma}
	\label{lem:FD1}
	Let $S$ be a solution for a given \textnormal{EMB} instance $I =(E,\cI,c,T)$,
	and let $P \in \cQ$ where $P = (\opt,\cF,\triangleleft)$. Then, $S$ is a solution for $R_P(I)$ of value $v_I(S) = k_I \cdot H_I+T$.  %be a non-trivial  \textnormal{MOL} problem, where $P = (\opt,\cF,\triangleleft)$. 
\end{lemma}

\begin{proof}
	Let $\cm = (E,\cI)$. Since $S$ is a solution for $I$ we have $S \in \bases(\cm)$; thus, $S \in \cF(\cm)$. Then,    
	$$w_{I,P}(S) = |S| \cdot H_I+c(S) \cdot (-1)^{d(P)} =  k_I \cdot H_I+T \cdot (-1)^{d(P)} = L_{I,P}.$$
	The second equality holds since $S$ is a solution for $I$; thus, $|S| = k_I$ (as $S$ is a basis of $\cm$), and $c(S) = T$. We conclude that $S$ is a solution for $R_P(I)$. Finally, note that $S$ satisfies 
%	\begin{equation*}
%		\label{eq:t=r}
%		v_I(S) = |S| \cdot H_I+c(S) =  k_I \cdot H_I+T.
%		\end{equation*}\end{proof}
\begin{align*}
~~~~~~~~~~~~~~~~~~~~~~~~~~~~~	v_I(S) = |S| \cdot H_I + c(S)
	= k_I \cdot H_I +T ~~~~~~~~~~~~~~~~~~~~~~~~~~~~~~~~~~~~~~~~~~~~~~~~~~~~~~~~ \qedhere
\end{align*}
\end{proof}

The next claim gives an upper bound on the optimal value for maximization MOL problems. We then derive an analogous lower bound for minimization (non-trivial) MOL problems. 











	\begin{lemma}
	\label{lem:OptMaxBound}
	Let $I =(E,\cI,c,T)$ be an \textnormal{EMB} instance and $P \in \cQ$, where $P = (\opt,\cF,\triangleleft)$ and $\left(\opt \textnormal{ is } \max\right)$. Then, for every solution $S$ of $R_P(I)$ it holds that $v_I(S) \leq k_I \cdot H_I+T$.  %be a non-trivial  \textnormal{MOL} problem, where $P = (\opt,\cF,\triangleleft)$. 
\end{lemma}

\begin{proof}
	Let $S$ be an optimal solution for $R_P(I)$. Thus, $|S| \leq k_I$ as $S \in \cI$. If $|S| < k_I$ then
	%\begin{itemize}
	%	\item If $d(P) = 0$. Then, consider the following sub cases.  
	 \begin{equation*}
			v_I(S) \leq (k_I-1) \cdot H_I+c(S) \leq (k_I-1) \cdot H_I+c(E) <k_I \cdot H_I \leq k_I \cdot H_I+T.
		\end{equation*}
	%	and 
%		\begin{equation}
%			\label{eq:S<Kv}
%			w_{I,P}(S) \leq (k-1) \cdot H_I+c(S)\cdot (-1)^{d(P)} \leq (k-1) \cdot H_I+c(E) <2k \cdot c(E) \leq k_I \cdot H_I+T.
%		\end{equation}
	Otherwise, $|S| = k_I$. Consider the two cases for $\triangleleft$.
			\begin{enumerate}
		\item $(\triangleleft \textnormal{ is}   \leq)$. Then, $d(P) = 0$ (see \Cref{def:reduction}); thus, since $S$ is a solution for $R_P(I)$: $$v_I(S) = w_{I,P}(S) \leq L_{I,P} = k_I \cdot H_I+T.$$ 
		\item $(\triangleleft \textnormal{ is}   \geq)$. Then, $d(P) = 1$. As $S$ is a solution for $R_P(I)$,
		\begin{equation}
			\label{eq:Dp=1}
			  k_I \cdot H_I-c(S)=|S|\cdot H_I-c(S) =w_{I,P}(S) \geq L_{I,P} = k_I \cdot H_I-T.
		\end{equation} 
		By \eqref{eq:Dp=1}, it follows that $c(S) \leq T$; thus, $	v_I(S) = k_I \cdot H_I+c(S) \leq k_I \cdot H_I+T$. 
	\end{enumerate} 
	% E.g., for $P = (\max,\IS,\leq)$ it holds that $d(P) = 0$ and for  $P' = (\max,\IS,\geq)$ it is $d(P') = 1$. 
	
%		\begin{enumerate}
%			\item $d(P) = 0$. Then, it follows that $\triangleleft = \leq $; thus, since $S$ is a solution for $R_P(I)$: $$v_I(S) = w_{I,P}(S) \leq L_{I,P} = k_I \cdot H_I+T.$$ 
%			\item $d(P) = 1$. Then, $\triangleleft = \geq $; thus, since $S$ is a solution for $R_P(I)$:
%			\begin{equation}
%				\label{eq:Dp=1}
%				w_{I,P}(S) = k \cdot H_I-c(S) \geq L_{I,P} = k \cdot H_I-T.
%			\end{equation} 
%			By \eqref{eq:Dp=1} it follows that $c(S) \leq T$; thus, $	v_I(S) = k \cdot H_I+c(S) \leq k_I \cdot H_I+T$. 
%		\end{enumerate} 
In all the above cases, we have that $v_I(S) \leq k_I \cdot H_I+T$, implying the statement of the lemma. 
\end{proof}

Now, for minimization problems we have the next result. 
	\begin{lemma}
	\label{lem:OptMinBound}
	Let $I =(E,\cI,c,T)$ be an \textnormal{EMB} instance, and $P \in \cQ$, where $P = (\opt,\cF,\triangleleft)$ and $\left(\opt \textnormal{ is } \min\right)$. Then, for every solution $S$ of $R_P(I)$, it holds that $v_I(S) \geq k_I \cdot H_I+T$.  %$\OPT\left( R_P(I) \right) \geq k_I \cdot H_I+T$.  %be a non-trivial  \textnormal{MOL} problem, where $P = (\opt,\cF,\triangleleft)$. 
\end{lemma}

\begin{proof}
	Let $S$ be a solution for $R_P(I)$. We consider two cases. 
	\begin{enumerate}
		\item $(\triangleleft \textnormal{ is}   \geq)$. Then, $d(P) = 0$; since $S$ is a solution for $R_P(I)$, $$v_I(S) = w_{I,P}(S) \geq L_{I,P} = k_I \cdot H_I+T.$$ 
		\item  $(\triangleleft \textnormal{ is}   \leq)$. Then, $d(P) = 1$; since $P$ is a non-trivial MOL problem (i.e., $P \in \cQ$), it follows that $\left(\cF \textnormal{ is } \bases\right)$. Thus, $S$ is a solution for $R_P(I)$ such that $|S| = k_I$, and
		\begin{equation}
			\label{eq:Dp=2}
			 	  k_I \cdot H_I-c(S)= |S| \cdot H_I-c(S) =w_{I,P}(S) \leq L_{I,P} = k_I \cdot H_I-T.
		\end{equation} 
		By \eqref{eq:Dp=2} it follows that $c(S) \geq T$; thus, $v_I(S) = k_I \cdot H_I+c(S) \geq k_I \cdot H_I+T$. 
	\end{enumerate}
	%	\end{itemize} 
%\item 
%\end{itemize}
In all of the above cases, we have that $v_I(S) \geq k_I \cdot H_I+T$, implying the statement of the lemma. 
\end{proof}


Using \Cref{lem:FD1,lem:FD2,lem:OptMaxBound,lem:OptMinBound}, we can now prove \Cref{thm:H}. 
\thmH*
\begin{proof}
%	The first property the theorem
We note that Property~\ref{h:1} follows directly from \Cref{lem:FD2}. For Property~\ref{h:2}, assume that $I$ is a ``yes"-instance, then by \Cref{lem:FD1}, there is a solution $D$ for $R_P(I)$ such that $v_I(D) = H_I \cdot k_I+T$. It remains to show Property~\ref{h:2}. $(ii)$.
%that every $(1+\eps_I)$-approximate solution $S$ for $R_P(I)$ satisfies $v_I(S) = H_I \cdot k_I+T$.
Let $S$ be a $(1+\eps_I)$-approximate solution for $R_P(I)$. We distinguish between two cases. 
	
		\begin{enumerate}
		\item $(\opt \textnormal{ is}   \max)$. Then, by \Cref{lem:OptMaxBound},
		\begin{equation*}
			\label{eq:Max1}
			0 \leq H_I \cdot k_I+T-v_I(S).
		\end{equation*} Moreover, since $S$ is a $(1+\eps_I)$-approximate solution for $R_P(I)$, and $D$ is a solution for $R_P(I)$,
		\begin{equation*}
			\label{eq:MaxApprox}
	%	0 \leq 
	H_I \cdot k_I+T-v_I(S) \leq H_I \cdot k_I+T - \frac{v_I(D)}{(1+\eps_I)} =  \frac{\eps_I \cdot \left(H_I \cdot k_I+T \right) }{(1+\eps_I)} \leq \eps_I \cdot \left(H_I \cdot k_I+T \right). 
		\end{equation*} %The first inequality follows from \Cref{lem:OptMaxBound}. 
	%	The second inequality holds since $S$ is a $(1+\eps_I)$-approximation for $R_P(I)$ and $D$ is some solution for $R_P(I)$. 
	By the above, it follows that%\eqref{eq:Max1} and \eqref{eq:MaxApprox} it follows that 
	\begin{equation*}
			  \left|	v_I(S)-\left(H_I \cdot k_I+T\right) \right|	\leq 	\eps_I \cdot \left(H_I \cdot k_I+T\right). 
	\end{equation*}
		\item  $(\opt \textnormal{ is}   \min)$. This case is analogous to the above.  By \Cref{lem:OptMinBound}, 
		\begin{equation*}
			\label{eq:Min1}
	0 \leq v_I(S) - \left(H_I \cdot k_I+T\right). 
		\end{equation*} Since $S$ is a $(1+\eps_I)$-approximate solution for $R_P(I)$, and $D$ is a solution for $R_P(I)$,
		\begin{equation*}
			\label{eq:MinApprox}
			%	0 \leq 
		v_I(S) - \left(H_I \cdot k_I+T\right) \leq (1+\eps_I) \cdot v_I(D) - \left(H_I \cdot k_I+T\right) =   \eps_I \cdot \left(H_I \cdot k_I+T \right). 
		\end{equation*} 
		By the above, 
		\begin{equation*}
			\left|	v_I(S)-\left(H_I \cdot k_I+T\right) \right|	\leq 	\eps_I \cdot \left(H_I \cdot k_I+T\right). 
		\end{equation*}

	\end{enumerate} 
	Thus in both cases it holds that, 
	\begin{equation}
		\label{eq:combined}
		\left|	v_I(S)-\left(H_I \cdot k_I+T\right) \right|	\leq 	\eps_I \cdot \left(H_I \cdot k_I+T\right). 
	\end{equation}
	Let $n = |E|+1$ and $m = c(E)+1$. Then, by the selection of $\eps_I$ in \eqref{eq:epsI}, 
	\begin{equation}
		\label{eq:MinMax}
		\eps_I \cdot \left(H_I \cdot k_I+T\right) = \frac{H_I \cdot k_I+T}{8 \cdot n \cdot (T+1) \cdot m} 
		\leq\frac{2 \cdot m \cdot n + T}{8 \cdot n \cdot (T+1) \cdot m} 
		<\frac{4 \cdot m \cdot n \cdot (T+1)}{8 \cdot n \cdot (T+1) \cdot m} 
		= \frac{1}{2}. 
	\end{equation}
The first inequality holds since $k_I \leq |E|$ and $H_I \leq 2 \cdot m$.	Therefore, by \eqref{eq:combined} and \eqref{eq:MinMax}, 
	\begin{equation}
		\label{eq:FinalMinMax}
		  \left|	v_I(S)-\left(H_I \cdot k_I+T\right) \right|	\leq 	\eps_I \cdot \left(H_I \cdot k_I+T\right) <  \frac{1}{2}.
	\end{equation} Since $v_I(S) \in \mathbb{N}$ by \Cref{def:reduction}, it follows from \eqref{eq:FinalMinMax} that $v_I(S) = H_I \cdot k_I+T$. This gives
	the statement of the theorem. 
\end{proof}

